% GENERATED FILE. Edit canonical lessons, then run npm run generate:books.
% canonical-source-hash: fnv1a64:5f2b2482da09d4fb
% canonical-lessons: MW-C270-maphi, MW-C270-kasur, MW-C270-pachhtano, MW-C270-phorno, MW-C270-khijno

\chapter{Saying Sorry}
\label{ch:mw-saying-sorry}
\hlchaptermodality{270}

\begin{chapteropening}
\textbf{By the end of this chapter:} \emph{I can say sorry and ask someone not to be cross.}

Complete the last lesson of chapter 270: I can say sorry and ask someone not to be cross.
\end{chapteropening}

\section[\texorpdfstring{\mw{माफी} (māphī)}{माफी (māphī)}]{\mw{माफी} (māphī) — pardon}

\label{lesson:MW-C270-maphi}



\begin{quote}
Before the new one: say the Marwadi for stubbornness, then the Marwadi for to interrupt.
\end{quote}

\subsection*{You'll want to know: \mw{माफी}}
\textbf{\mw{माफी}} — \emph{māphī} — \textquotedblleft{}pardon\textquotedblright{}.

\begin{quote}
\mw{हूं} \mw{माफी} \mw{मांगूं} \mw{हूं।}
\end{quote}

\begin{quote}
\emph{hū̃ māphī māṅgū̃ hū̃.}
\end{quote}

\begin{quote}
I'm sorry: I ask your pardon.
\end{quote}

A feminine noun.

\begin{cousinweb}
\textbf{\mw{माफी}} (\emph{māphī}) is usually traced, through Persian \emph{muʿāfī}, \textquotedblleft{}pardon\textquotedblright{}, to Arabic \emph{muʿāf}, \textquotedblleft{}excused\textquotedblright{}.
\end{cousinweb}

\begin{grammarlens}[title={What you've built}]
One more word for asking and answering.
\end{grammarlens}

\subsection*{Your turn}
\begin{itemize}
\raggedright
  \item \emph{Say it:} \emph{māphī}
  \item \emph{Say it:} \emph{māphī}, once more
  \item \emph{Say it:} \emph{hū̃ māphī māṅgū̃ hū̃.} — I'm sorry: I ask your pardon.
  \item \emph{Recall:} say the Marwadi for stubbornness, then the Marwadi for to interrupt, then say \emph{māphī} again
\end{itemize}

\begin{tcolorbox}[breakable,colback=teal!4,colframe=teal!35!black,title={Before you move on}]
What does \mw{माफी} mean? (\textquotedblleft{}Pardon\textquotedblright{}.) Say it once more.
\end{tcolorbox}
\section[\texorpdfstring{\mw{कसूर} (kasūr)}{कसूर (kasūr)}]{\mw{कसूर} (kasūr) — a wrongdoing; a fault}

\label{lesson:MW-C270-kasur}



\begin{quote}
Before the new one: say the Marwadi for to interrupt, then the Marwadi for pardon.
\end{quote}

\subsection*{You'll want to know: \mw{कसूर}}
\textbf{\mw{कसूर}} — \emph{kasūr} — \textquotedblleft{}a wrongdoing; a fault\textquotedblright{}.

\begin{quote}
\mw{थारो} \mw{कोई} \mw{कसूर} \mw{कोनी।}
\end{quote}

\begin{quote}
\emph{thāro koī kasūr konī.}
\end{quote}

\begin{quote}
It isn't your fault at all.
\end{quote}

A masculine noun.

\begin{cousinweb}
\textbf{\mw{कसूर}} (\emph{kasūr}) is usually traced, through Persian \emph{qusūr}, to Arabic \emph{quṣūr}, \textquotedblleft{}falling short\textquotedblright{}.
\end{cousinweb}

\begin{grammarlens}[title={What you've built}]
One more word for asking and answering.
\end{grammarlens}

\subsection*{Your turn}
\begin{itemize}
\raggedright
  \item \emph{Say it:} \emph{kasūr}
  \item \emph{Say it:} \emph{kasūr}, once more
  \item \emph{Say it:} \emph{thāro koī kasūr konī.} — It isn't your fault at all.
  \item \emph{Recall:} say the Marwadi for to interrupt, then the Marwadi for pardon, then say \emph{kasūr} again
\end{itemize}

\begin{tcolorbox}[breakable,colback=teal!4,colframe=teal!35!black,title={Before you move on}]
What does \mw{कसूर} mean? (\textquotedblleft{}A wrongdoing; a fault\textquotedblright{}.) Say it once more.
\end{tcolorbox}
\section[\texorpdfstring{\mw{पछताणो} (pachhtāṇo)}{पछताणो (pachhtāṇo)}]{\mw{पछताणो} (pachhtāṇo) — to regret}

\label{lesson:MW-C270-pachhtano}



\begin{quote}
Before the new one: say the Marwadi for pardon, then the Marwadi for a wrongdoing.
\end{quote}

\subsection*{You'll want to know: \mw{पछताणो}}
\textbf{\mw{पछताणो}} — \emph{pachhtāṇo} — \textquotedblleft{}to regret\textquotedblright{}.

\begin{quote}
\mw{हूं} \mw{घणो} \mw{पछताऊं} \mw{हूं।}
\end{quote}

\begin{quote}
\emph{hū̃ ghaṇo pachhtāū̃ hū̃.}
\end{quote}

\begin{quote}
I am very sorry for it.
\end{quote}

Past (he …): \mw{वो} \mw{पछतायो}, \emph{vo pachhtāyo}.

\begin{cousinweb}
\textbf{\mw{पछताणो}} (\emph{pachhtāṇo}) is inherited from a reconstructed \emph{*paścottāpa}, \textquotedblleft{}repentance\textquotedblright{}; CDIAL lists Old Marwari \emph{pichatāi} and \emph{pachitāi} under it (CDIAL 8010).
\end{cousinweb}

\begin{grammarlens}[title={What you've built}]
One more word for asking and answering.
\end{grammarlens}

\subsection*{Your turn}
\begin{itemize}
\raggedright
  \item \emph{Say it:} \emph{pachhtāṇo}
  \item \emph{Say it:} \emph{pachhtāṇo}, once more
  \item \emph{Say it:} \emph{hū̃ ghaṇo pachhtāū̃ hū̃.} — I am very sorry for it.
  \item \emph{Recall:} say the Marwadi for pardon, then the Marwadi for a wrongdoing, then say \emph{pachhtāṇo} again
\end{itemize}

\begin{tcolorbox}[breakable,colback=teal!4,colframe=teal!35!black,title={Before you move on}]
What does \mw{पछताणो} mean? (\textquotedblleft{}To regret\textquotedblright{}.) Say it once more.
\end{tcolorbox}
\section[\texorpdfstring{\mw{फोड़णो} (phoṛṇo)}{फोड़णो (phoṛṇo)}]{\mw{फोड़णो} (phoṛṇo) — to smash; to break open}

\label{lesson:MW-C270-phorno}



\begin{quote}
Before the new one: say the Marwadi for a wrongdoing, then the Marwadi for to regret.
\end{quote}

\subsection*{You'll want to know: \mw{फोड़णो}}
\textbf{\mw{फोड़णो}} — \emph{phoṛṇo} — \textquotedblleft{}to smash; to break open\textquotedblright{}.

\begin{quote}
\mw{मटको} \mw{मत} \mw{फोड़ो।}
\end{quote}

\begin{quote}
\emph{maṭko mat phoṛo.}
\end{quote}

\begin{quote}
Don't smash the pot.
\end{quote}

Past (I …): \mw{म्हैं} \mw{मटको} \mw{फोड़्यो}, \emph{mhaĩ maṭko phoṛyo}.

\begin{cousinweb}
\textbf{\mw{फोड़णो}} (\emph{phoṛṇo}) is inherited from Sanskrit \emph{sphoṭayati}, \textquotedblleft{}bursts, splits, breaks\textquotedblright{}; CDIAL gives no Marwari form in this entry, but lists Hindi \emph{phoṛnā}, Gujarati \emph{phoṛvũ} and Sindhi \emph{phoṛaṇu}, all \textquotedblleft{}to break\textquotedblright{} (CDIAL 13857). It is the transitive partner of \textbf{\mw{फूटणो}}, to crack open: \textbf{\mw{मटको} \mw{फूट्यो}}, but \textbf{\mw{म्हैं} \mw{मटको} \mw{फोड़्यो}}.
\end{cousinweb}

\begin{grammarlens}[title={What you've built}]
One more word for asking and answering.
\end{grammarlens}

\subsection*{Your turn}
\begin{itemize}
\raggedright
  \item \emph{Say it:} \emph{phoṛṇo}
  \item \emph{Say it:} \emph{phoṛṇo}, once more
  \item \emph{Say it:} \emph{maṭko mat phoṛo.} — Don't smash the pot.
  \item \emph{Recall:} say the Marwadi for a wrongdoing, then the Marwadi for to regret, then say \emph{phoṛṇo} again
\end{itemize}

\begin{tcolorbox}[breakable,colback=teal!4,colframe=teal!35!black,title={Before you move on}]
What does \mw{फोड़णो} mean? (\textquotedblleft{}To smash; to break open\textquotedblright{}.) Say it once more.
\end{tcolorbox}
\section[\texorpdfstring{\mw{खीजणो} (khījṇo)}{खीजणो (khījṇo)}]{\mw{खीजणो} (khījṇo) — to get cross}

\label{lesson:MW-C270-khijno}



\begin{quote}
Before the new one: say the Marwadi for to regret, then the Marwadi for to smash.
\end{quote}

\subsection*{You'll want to know: \mw{खीजणो}}
\textbf{\mw{खीजणो}} — \emph{khījṇo} — \textquotedblleft{}to get cross\textquotedblright{}.

\begin{quote}
\mw{मत} \mw{खीजो} \mw{सा।}
\end{quote}

\begin{quote}
\emph{mat khījo sā.}
\end{quote}

\begin{quote}
Please don't be cross.
\end{quote}

Past (he …): \mw{वो} \mw{खीज्यो}, \emph{vo khījyo}.

\begin{cousinweb}
\textbf{\mw{खीजणो}} (\emph{khījṇo}) is inherited from Sanskrit \emph{kṣīyate}, \textquotedblleft{}is worn away, is injured\textquotedblright{}; CDIAL lists Old Marwari \emph{khijaï}, \textquotedblleft{}to be irritated\textquotedblright{}, and Gujarati \emph{khījvũ}, \textquotedblleft{}to take offence\textquotedblright{} (CDIAL 3695).
\end{cousinweb}

\begin{grammarlens}[title={What you've built}]
One more word for asking and answering.
\end{grammarlens}

\subsection*{Your turn}
\begin{itemize}
\raggedright
  \item \emph{Say it:} \emph{khījṇo}
  \item \emph{Say it:} \emph{khījṇo}, once more
  \item \emph{Say it:} \emph{mat khījo sā.} — Please don't be cross.
  \item \emph{Recall:} say the Marwadi for to regret, then the Marwadi for to smash, then say \emph{khījṇo} again
\end{itemize}

\begin{tcolorbox}[breakable,colback=teal!4,colframe=teal!35!black,title={Before you move on}]
What does \mw{खीजणो} mean? (\textquotedblleft{}To get cross\textquotedblright{}.) Say it once more.
\end{tcolorbox}
